% TOP_PROBLEM: {"id":30007742,"problem_number":"LOCAL-30007742","title":"Housing-supply responses to reform","category_id":21,"category":{"id":21,"name":"economics","display_name":"Economics","description":"Open economic theory statements, published research agendas, and proposed scoped specifications; question provenance and historical priority are qualified per record.","slug":"economics","order_index":21,"created_at":"2026-10-08T00:00:00Z"},"status":"research_question","external_url":null,"legacy_ids":[],"published":false,"statement":"Estimate a fifteen-year profile of housing-stock and quality-adjusted rent responses to a specified density reform, including delays and neighboring-area spillovers.","economics":{"original_id":231,"field":"Urban, housing and spatial economics","kind":"empirical","formalization_basis":"adapted_research_question","source_status":"explicit_open","priority_status":"earliest_found","priority_note":"Earliest explicit statement located in this audit; no exhaustive priority search or first-ever claim. The mathematical specification below is adapted, not quoted from the source.","earliest_verified_reference":{"authors":["Nathaniel Baum-Snow","Gilles Duranton"],"title":"Housing supply and housing affordability","year":2025,"url":"https://realestate.wharton.upenn.edu/wp-content/uploads/2025/12/BaumSnow_Duranton_bc_2025-003.pdf","doi":"10.1016/bs.hesreg.2025.05.003","locator":"§7 Conclusions and future research directions, printed pp. 448–449","evidence_summary":"The authors explicitly seek dynamic supply elasticities, construction-productivity causes and remedies, and richer spatial models of migration and housing reform."},"current_reference":{"authors":["Nathaniel Baum-Snow","Gilles Duranton"],"title":"Housing supply and housing affordability","year":2025,"url":"https://realestate.wharton.upenn.edu/wp-content/uploads/2025/12/BaumSnow_Duranton_bc_2025-003.pdf","doi":"10.1016/bs.hesreg.2025.05.003","locator":"§7 Conclusions and future research directions, printed pp. 448–449","evidence_summary":"The authors explicitly seek dynamic supply elasticities, construction-productivity causes and remedies, and richer spatial models of migration and housing reform."},"formalization_note":"The source explicitly identifies the underlying gap or agenda. Population, horizon, interventions, estimands, and auxiliary assumptions are proposed operational choices, not a canonical source theorem. Partial later answers are compatible with this scoped question; the citation alone does not prove absence of every subsequent solution. Mathematical scope was checked independently; added definedness, feasibility, or identification conditions belong to this proposed formalization."},"statusQualification":"Published question or research agenda verified at the cited locator. The mathematical specification is adapted for this catalogue and includes additional assumptions; source publication does not establish first-ever priority or certify this exact model remains unsolved. The source explicitly identifies the underlying gap or agenda. Population, horizon, interventions, estimands, and auxiliary assumptions are proposed operational choices, not a canonical source theorem. Partial later answers are compatible with this scoped question; the citation alone does not prove absence of every subsequent solution. Mathematical scope was checked independently; added definedness, feasibility, or identification conditions belong to this proposed formalization.","statusReviewedAt":"2026-10-08"}
% FIRST_POSTED: 2026-10-08
% ENTRY_KIND: statement_only
% !TeX program = lualatex
\documentclass[11pt]{article}
\usepackage{fontspec}
\usepackage[a4paper,margin=25mm]{geometry}
\usepackage{amsmath,amssymb,mathtools}
\usepackage{xurl}
\usepackage[hidelinks]{hyperref}
\newcommand{\sref}[2]{\hyperlink{source:#1:#2}{[S#2]}}
\setlength{\emergencystretch}{3em}
\begin{document}
% problemId:30007742
\section*{Housing-supply responses to reform}\label{30007742}
\subsection{Definitions and mathematical statement}
Consider metropolitan areas in one country adopting a specified reform that increases permitted residential density. Let \(a=1\) denote implementation and \(a=0\) continuation of prior rules. Define \(H_{mt}(a)\) as habitable housing units, \(R_{mt}(a)\) quality-adjusted real rent, and \(t=0\) as the reform year in metropolitan area \(m\). For every \(h\in\{1,\ldots,15\}\), define
\[\tau_H(h)=E[\log H_{m,h}(1)-\log H_{m,h}(0)],\qquad\tau_R(h)=E[\log R_{m,h}(1)-\log R_{m,h}(0)].\]
Require positive housing stocks and rents under both policies and integrable log contrasts. Identify these dynamic responses using randomized rollout or a justified comparison design, allowing construction delays, demolition, redevelopment, and infrastructure constraints. Record actual permitting and enforcement, since statutory eligibility does not equal realized construction. Keep baseline metropolitan boundaries and report spillovers to adjoining areas. An estimate at one horizon does not identify a complete supply elasticity: separately specify demand shocks and any price-based elasticity target. Determine which local conditions alter response speed and transport estimates only to areas with adequate covariate overlap. The housing review explicitly calls for systematic dynamic market-level supply evidence; this particular reform, treatment counterfactual, and fifteen-year profile are an adapted operationalization.

\par\textbf{Formulation and scope.} The source explicitly identifies the underlying gap or agenda. Population, horizon, interventions, estimands, and auxiliary assumptions are proposed operational choices, not a canonical source theorem. Partial later answers are compatible with this scoped question; the citation alone does not prove absence of every subsequent solution. Mathematical scope was checked independently; added definedness, feasibility, or identification conditions belong to this proposed formalization.

\par\textbf{Open-question provenance.} An explicit question or research agenda was verified at the source locator. This does not by itself certify that every later version remains unresolved \sref{30007742}{1}.

\par\textbf{Historical priority.} Earliest explicit statement located in this audit; no exhaustive priority search or first-ever claim. The mathematical specification below is adapted, not quoted from the source.

\subsection{Short English statement}
Estimate a fifteen-year profile of housing-stock and quality-adjusted rent responses to a specified density reform, including delays and neighboring-area spillovers.

\subsection{Sources}
\begin{itemize}\raggedright
\item[\textnormal{[S1]}]\hypertarget{source:30007742:1}{} Nathaniel Baum-Snow, Gilles Duranton. 2025. Housing supply and housing affordability. §7 Conclusions and future research directions, printed pp. 448–449.
\url{https://realestate.wharton.upenn.edu/wp-content/uploads/2025/12/BaumSnow_Duranton_bc_2025-003.pdf}.
DOI: 10.1016/bs.hesreg.2025.05.003.
{\small\textit{Earliest verified question/agenda reference; historical priority qualified below. The authors explicitly seek dynamic supply elasticities, construction-productivity causes and remedies, and richer spatial models of migration and housing reform.}}
\end{itemize}
\end{document}
